% GENERATED FILE. Edit canonical lessons, then run npm run generate:books.
% canonical-source-hash: fnv1a64:1fa7535add818308
% canonical-lessons: ML-C53-sky, ML-C53-sun, ML-C53-moon, ML-C53-star, ML-C53-sunshine

\chapter{Overhead}
\label{ch:ml-overhead}
\hlchaptermodality{53}

\begin{chapteropening}
\textbf{By the end of this chapter:} \emph{I can name the sky and four things in it in Malayalam.}

Complete ML-C53-sunshine and say all five words in order.
\end{chapteropening}

\section[ākāśaṁ]{\ml{ആകാശം} (ākāśaṁ) — the sky}

\label{lesson:ML-C53-sky}



\begin{quote}
Before the new run of words: what did \emph{pūvŭ} mean, and what did \emph{māla} mean?
\end{quote}

\subsection*{You'll want to know: \ml{ആകാശം}}
\textbf{\ml{ആകാശം}} (\emph{ākāśaṁ}) — \textquotedblleft{}the sky\textquotedblright{}.

Sanskrit \emph{ākāśa}, built on a piece meaning \textquotedblleft{}to shine, to be visible\textquotedblright{}. The sky is named there for being the open lit place where a thing can be seen at all, rather than for being a lid over the world.

Malayalam keeps an inherited word for the same overhead space, \ml{വാനം} (\emph{vānaṁ}), which owes Sanskrit nothing. The two divide the work between them: \ml{ആകാശം} takes the ordinary shift and turns up in weather talk, while \ml{വാനം} keeps to songs and to verse. A borrowed word taking the daily job and leaving the older one the ceremony is an arrangement this book has shown you before, at the doorway.

\emph{ākāśa} travelled as well. Traders carried it east, and Malay and Indonesian still say \emph{angkasa} for the sky.

\begin{grammarlens}[title={What you've built}]
The space the rest of this run hangs in.
\end{grammarlens}

\subsection*{Your turn}
Say these aloud:

\begin{itemize}
\raggedright
  \item \emph{ākāśaṁ}
  \item it once more, slowly
  \item \emph{ākāśaṁ}, then \emph{namaskāraṁ}, as you would greet somebody out of doors
\end{itemize}

\begin{tcolorbox}[breakable,colback=teal!4,colframe=teal!35!black,title={Before you move on}]
What does \ml{ആകാശം} mean? (\textquotedblleft{}the sky\textquotedblright{}.) Say it once more.
\end{tcolorbox}
\section[sūryan]{\ml{സൂര്യൻ} (sūryan) — the sun}

\label{lesson:ML-C53-sun}



\begin{quote}
Before the new one: what did \emph{ākāśaṁ} mean?
\end{quote}

\subsection*{You'll want to know: \ml{സൂര്യൻ}}
\textbf{\ml{സൂര്യൻ}} (\emph{sūryan}) — \textquotedblleft{}the sun\textquotedblright{}.

Sanskrit \emph{sūrya}, wearing the Malayalam chillu \ml{ൻ} on the end — the same tail that closes \emph{adhyāpakan} and \emph{karṣakan}, and it marks this one out as a he-word. The sun is a person in that grammar, not a thing.

The Sanskrit word has cousins a long way west. It shares an ancestor with Latin \emph{sol} and with the Greek piece behind \emph{helios}, which is why English \emph{solar} and \emph{heliotrope} are distant relatives of the word a Malayali newsreader says.

Malayalam's own sun-word is \ml{ഞായർ} (\emph{ñāyaṟ}), and it survives mostly inside the name of Sunday, \ml{ഞായറാഴ്ച} (\emph{ñāyaṟāḻca}) — the sun's day, arrived at independently of the English one.

\begin{grammarlens}[title={What you've built}]
Two: the sky, and the thing that crosses it.
\end{grammarlens}

\subsection*{Your turn}
Say these aloud:

\begin{itemize}
\raggedright
  \item \emph{sūryan}
  \item it once more, slowly
  \item \emph{ākāśaṁ}, then \emph{sūryan}, and say which of the two is the container
\end{itemize}

\begin{tcolorbox}[breakable,colback=teal!4,colframe=teal!35!black,title={Before you move on}]
What does \ml{സൂര്യൻ} mean? (\textquotedblleft{}the sun\textquotedblright{}.) Say it once more.
\end{tcolorbox}
\section[candran]{\ml{ചന്ദ്രൻ} (candran) — the moon}

\label{lesson:ML-C53-moon}



\begin{quote}
Before the new one: what did \emph{sūryan} mean?
\end{quote}

\subsection*{You'll want to know: \ml{ചന്ദ്രൻ}}
\textbf{\ml{ചന്ദ്രൻ}} (\emph{candran}) — \textquotedblleft{}the moon\textquotedblright{}.

Sanskrit \emph{candra} means \textquotedblleft{}shining, bright\textquotedblright{}, and Malayalam took it whole, then gave it the same he-word tail the sun got. Sun and moon travel as a pair in this language, down to the shape of their endings.

\emph{candra} went west too. It shares an ancestor with the Latin piece behind English \emph{candle} and \emph{incandescent}: a very old root for burning white.

Malayalam's inherited moon-word is \ml{തിങ്കൾ} (\emph{tiṅkaḷ}), and that word also names Monday and once named a month. One shape for the moon and for the stretch of time it measures — the moon being the oldest calendar there is, and Malayalam keeping that fact inside a single word.

\begin{grammarlens}[title={What you've built}]
Three, and the last two were borrowed together.
\end{grammarlens}

\subsection*{Your turn}
Say these aloud:

\begin{itemize}
\raggedright
  \item \emph{candran}
  \item it once more, slowly
  \item \emph{sūryan}, then \emph{candran}, and hear the same tail on both
\end{itemize}

\begin{tcolorbox}[breakable,colback=teal!4,colframe=teal!35!black,title={Before you move on}]
What does \ml{ചന്ദ്രൻ} mean? (\textquotedblleft{}the moon\textquotedblright{}.) Say it once more.
\end{tcolorbox}
\section[nakṣatraṁ]{\ml{നക്ഷത്രം} (nakṣatraṁ) — a star}

\label{lesson:ML-C53-star}



\begin{quote}
Before the new one: what did \emph{candran} mean?
\end{quote}

\subsection*{You'll want to know: \ml{നക്ഷത്രം}}
\textbf{\ml{നക്ഷത്രം}} (\emph{nakṣatraṁ}) — \textquotedblleft{}a star\textquotedblright{}.

Sanskrit \emph{nakṣatra}. In Kerala it does a job it does not do in English: a person's \ml{നക്ഷത്രം} is the birth-star, one of twenty-seven asterisms the moon passes through, and it is the star you were born under rather than the date that fixes when your birthday falls.

So a Malayali may keep two birthdays a year apart from each other — the calendar one and the star one — and the star one is the one the household cooks for.

Tamil reaches for the same borrowed word, \emph{naṭcattiram}, worn down a little further. The inherited Dravidian way of naming a star was to call it a fish in the sky, and Tamil still has \emph{viṇmīn} for it; Malayalam let that one go.

\begin{grammarlens}[title={What you've built}]
Four, and this one keeps a calendar.
\end{grammarlens}

\subsection*{Your turn}
Say these aloud:

\begin{itemize}
\raggedright
  \item \emph{nakṣatraṁ}
  \item it once more, slowly
  \item \emph{nakṣatraṁ}, then \emph{candran}, and say which of the two moves through the other
\end{itemize}

\begin{tcolorbox}[breakable,colback=teal!4,colframe=teal!35!black,title={Before you move on}]
What does \ml{നക്ഷത്രം} mean? (\textquotedblleft{}a star\textquotedblright{}.) Say it once more.
\end{tcolorbox}
\section[veyil]{\ml{വെയിൽ} (veyil) — sunshine, the sun's heat}

\label{lesson:ML-C53-sunshine}



\begin{quote}
Before the new one: what did \emph{nakṣatraṁ} mean?
\end{quote}

\subsection*{You'll want to know: \ml{വെയിൽ}}
\textbf{\ml{വെയിൽ}} (\emph{veyil}) — \textquotedblleft{}sunshine, the sun's heat\textquotedblright{}.

After four borrowed words, an inherited one. \ml{വെയിൽ} is Malayalam's own, matching Tamil \emph{veyil} exactly, and it names the heat that falls out of the sky rather than the body that sends it.

That is worth pausing on. English talks about the weather by naming the sun; Malayalam reaches first for what the sun does to you standing in it. A Kochi sentence about the afternoon is far more likely to be about \ml{വെയിൽ} than about \emph{sūryan} at all. Step out of the shade and you have stepped into \ml{വെയിൽ}.

Beside it sits its opposite, \ml{തണൽ} (\emph{taṇal}), the shade — and between those two words most of a hot afternoon can be described without the sun being mentioned once.

\begin{grammarlens}[title={What you've built}]
Five: the sky, the sun, the moon, a star, and what comes down out of it by day.
\end{grammarlens}

\subsection*{Your turn}
Say these aloud:

\begin{itemize}
\raggedright
  \item \emph{veyil}
  \item it once more, slowly
  \item all five in order — \emph{ākāśaṁ}, \emph{sūryan}, \emph{candran}, \emph{nakṣatraṁ}, \emph{veyil}
\end{itemize}

\begin{tcolorbox}[breakable,colback=teal!4,colframe=teal!35!black,title={Before you move on}]
What does \ml{വെയിൽ} mean? (\textquotedblleft{}sunshine, the sun's heat\textquotedblright{}.) Say it once more.
\end{tcolorbox}
